
\begin{figure*}
    \centering
    \includegraphics[width=\linewidth]{figures/overview2.png}
    \caption{Overview (Draft)}
    \label{fig:overview}
\end{figure*}
\section{Methods}



We now introduce the proposed multimodal localization method. Our method consists of an innovative WiFi localization neural network and a bespoken Bayesian filtering approach fusing together the WiFi localization results, odometry estimates and noisy radio map.
Fig.~\ref{pipeline} illustrates our proposed framework \chris{@Josh, we need a figure here to explain things}. 



\subsection{WiFi Fingerprint Transformer Network Design}

As found in \cite{wifioverview2017xia} 
% \chris{@Josh, we need a reference here.}
, the lengths of WiFi fingerprints $\mathbf{W}_i$ vary across locations and time even at the same location. This requires the WiFi localization model to take WiFi fingerprints in different lengths but still be able to output location coordinates. This requirement is similar to the one of natural language processing, where language models often take input as the sentences in different token lengths, but are still able to perform inference. Inspired by recent successes of transformers in natural language processing \cite{vaswani2017attention,devlin2018bert}, 
we propose utilizing the transformer's encoder in our network design to handle the varying-length WiFi fingerprints.
% we propose to design a transformer-like neural network architecture to process the varying-length WiFi fingerprints. 



\begin{figure}
    \centering
    \includegraphics[width=\linewidth]{figures/createRSSmap.png}
    \caption{Extract RSS Distribution from Radio Map}
    \label{fig:createrssi}
\end{figure}

\subsubsection{WiFi Fingerprint Embedding}
A fingerprint in the list format of $(MAC_j, RSS_j)$ cannot be directly employed as the input to a transformer-like network. 
Akin to the idea of distributed representations of words in natural language processing \cite{mikolov2013efficient}, 
a fingerprint need to be represented as a continuous vector before being fed into the model. 
Pre-trained word embedding layers are often available for language models by which the derived word embeddings are close to each other if two words have similar semantic meanings. 
However, such a semantic-preserving embedding layer is absent in the case of WiFi fingerprints esp. for the MAC address of APs. The challenge lies in that it is difficult to relate two nearby APs by their MAC address manifestation - which is simply a 12-character alphanumeric attribute that is used to identify individual electronic devices. To inject the location vicinity semantics in MAC address embedding, we propose to utilize the RSS distribution map (see Fig.\ref{fig:createrssi}) associated with each WiFi MAC address. The pixel value of a distribution map corresponds to the WiFi RSS value index at a particular location retrieved from the radio map. 
Intuitively, the MAC addresses of nearby WiFi APs tend to have similar RSS distribution (see Fig.\ref{}). 
The distinctions between APs can be robustly represented within the RSS distribution map, serving as a priori knowledge.
We further use a convolutional neural network to process the distribution maps and extract the MAC address embeddings into continuous vectors. 
During training, the physical vicinity can be conveyed in the similarity of the embeddings. Even for the same AP, its RSS values measured by a mobile agent can vary across places. 
To make use of the RSS value in the WiFi fingerprint,  we further introduce a learnable RSS encoding layer that maps the RSS value  between 0 and 1.  This scalar is multiplied with the MAC address embedding to derive the final fingerprint embedding, $\mathbf{E}_i$.  

% Analogous to the approach in transformer models\cite{vaswani2017attention},   where positional encoding vectors are added to each embedding to incorporate positional information, we have multiplied each MAC address embedding with a learnable scalar to represent its RSS value.
% Given that the relationships among RSS values exhibit nonlinear characteristics, which can vary significantly depending on the architectural attributes\cite{wifioverview2017xia}, an additional neural network, $\sigma$ in Fig.\ref{fig:AP2token}, has been integrated as a positional encoding to map RSS values to a scalar ranging between 0 and 1. After multiplication, the MAC has become the input MAC-RSS token that contains information about the RSS value.

\subsubsection{Multi-Head Attention Block.}

We aim to accurately predict location coordinates based on WiFi fingerprint embeddings. Given the dynamic and complex nature of WiFi fingerprints, it's crucial to capture dependencies between different features of the fingerprint for a precise location prediction. The self-attention transformer offers a promising solution, as it excels in modelling such dependencies in data. 

We thus adopt the self-attention transformer. Essentially, it allows each feature in the WiFi fingerprint embedding to weigh the importance of other features based on their content. This is particularly advantageous for our task because certain WiFi fingerprints might carry more relevance in determining specific location coordinates than others. Diving deeper into the mechanics, let's first revisit the single-head attention mechanism:
\begin{equation}
\begin{aligned}
head_h = Attention(\mathbf{E}_i \mathbf{W}_{Q_h}, \mathbf{E}_i \mathbf{W}_{K_h}, \mathbf{E}_i \mathbf{W}_{V_h}) \\
= softmax\left(\frac{\mathbf{E}_i \mathbf{W}_{Q_h}(\mathbf{E}_i \mathbf{W}_{K_h})^T}{\sqrt{d_k}}\right)\mathbf{E}_i \mathbf{W}_{V_h}
\end{aligned}
\end{equation}
where $E$ is the input fingerprint embedding; $W_{Qi}$, $W_{Ki}$, and $W_{Vi}$ are the projection matrices for query, key, and value respectively; $\sqrt{d_k}$ is the dimension of the key $k$ and query $q$.

However, relying on a single attention head might restrict the model from capturing multiple types of relationships or dependencies between the features. This is where the multi-head attention mechanism, as proposed in \cite{vaswani2017attention}, comes into play:
\begin{equation}
\begin{aligned}
MultiHead(\mathbf{E}_i \mathbf{W}_{Q_h}, \mathbf{E}_i \mathbf{W}_{K_h}, \mathbf{E}_i \mathbf{W}_{V_h})\\=Concat(head_1, \ldots, head_H) \cdot \mathbf{W}_O
\end{aligned}
\end{equation}
where $head_1, \ldots, head_H$ represent the $H$ single-head attention results. By using multiple attention heads, the model can simultaneously attend to various subspace representations of the input. Each head potentially captures different types of dependencies within the fingerprint embeddings. 
These concatenated results are then combined using a linear transformation matrix $W_O$, to produce an attended embedding. Each attended embedding is applied a non-linear transformation before output. After processing through several of such multi-head attention layers, the output from the last layer is further refined by a feed-forward network to regress the location coordinates.


\subsubsection{Uncertainty Estimation} 

As mentioned earlier, crowdsourced radio maps often come with inherent noise. Using them directly to train a transformer localization model can lead to high susceptibility to errors in predictions, given the model might learn from the noise rather than the actual signal. As a natural response to this challenge, one might consider ways to make the model more resilient to noise and simultaneously offer a more informative prediction. Thus, instead of merely giving a point estimate which provides a singular location prediction, it would be more beneficial if our model could quantify the level of confidence in its predictions. 

To this end, we reformulate the network's output of location coordinates as a \emph{probability distribution}. We choose the Gaussian distribution for this purpose, given its well-understood properties and ease of integration in subsequent fusion steps:
\begin{equation}
\mathbf{\check{L}}_i \sim \mathcal{N}(f_{\mu}(\mathbf{G}_i), f_{\sigma}(\mathbf{G}_i))
\end{equation}
% \begin{equation}
% \mathbf{\check{L}}_i \sim \mathcal{N}(f_{\mu}(\mathbf{W}_i), f_{\sigma}(\mathbf{W}_i))
% \end{equation}
where $f_{\sigma}(\cdot)$ is an extra neural network head capturing the variance which is pivotal as it provides an estimate of the uncertainty associated with the prediction. By transitioning to this probabilistic framework, we ensure the transformer model does not blindly commit to a potentially noisy sample. Instead, it offers a range of possibilities, weighted by their likelihood. This not only makes the model less prone to overfitting to the noise but also enriches the output, as the variance gives an insight into the reliability of the prediction. 

Altogether, given crowdsourced a radio map $\mathcal{M} = \{(\mathbf{W}_i, \mathbf{L}_i)|i=1,2,..,N\}$, we train our WiFi localization neural network with an uncertainty-aware loss function:
\begin{equation}
\mathcal{L} = \sum_{i}^{N} \frac{\lVert f_{\mu}(\mathbf{G}_i) - \mathbf{L}_i \rVert_2}{f_{\sigma}(\mathbf{G}_i)}+ln{f_{\sigma}(\mathbf{G}_i)} 
\end{equation}
% \begin{equation}
% \mathcal{L} = \sum_{i}^{N} \frac{\lVert f_{\mu}(\mathbf{W}_i) - \mathbf{L}_i \rVert_2}{f_{\sigma}(\mathbf{W}_i)}+ln{f_{\sigma}(\mathbf{W}_i)}
% \end{equation}

During training, the first term drives predictions closer to the radio map location coordinates. The numerator of the first term is to ensure that a larger variance is given for samples with significant prediction errors, mitigating large gradient impacts. The second term acts as a regularizer, discouraging the network from consistently yielding high variance.


\subsection{Multimodal Fusion via extended Kalman particle filter}
The extended Kalman filter is an improved Kalman filter for nonlinear systems, which
may provide unreliable estimates when the system is nonlinear. 
Particle filters are a predominant method to fuse multimodal data for robust localization. Prior arts have shown the effectiveness of using particle filters to combine the WiFi localization results, inertial odometry and floor plans. However, whether crowdsourced radio maps can serve as an alternative to floor plans in fusion remains largely unknown. Given a raw radio map, the kernel density estimation algorithm\cite{} can be used to generate a continuous likelihood distribution over the map area. However, the continuous likelihood distribution is imbalanced and noisy because the crowdsourced radio map itself is spares and noisy. Because of the absence of particles in true areas during the filter convergence\cite{Bukhori2017Detection}, the kidnapped robot problem \cite{Thrun2005Probabilistic,hester2008negative} is likely to happen and incurs localization failures. Therefore we combine the standard PF with the EKF to overcome this drawback of the PF. \iffalse Compared to the regular EKF, this method increases computational costs but can achieve higher accuracy and makes it possible to use a crowdsourced radio map to constrain the positioning results.\fi The Kalman correction can modify the position of the particle directly based on the uncertainty-aware WiFi localization result, to help the particles recover. \chris{@Xiangyu, I somehow get confused here. Why the normal prediction step in PF cannot help the particles recover?}

\subsubsection{prediction stage} 
In the extended Kalman particle filter, every particle maintains states of the particle $\mathbf{s}_{t}=[x_t,y_t,\theta_t,v_t]^T$, an importance weight $w$ and a corresponding error covariance matrix $\mathbf{P}_{t}$, where $x_t,y_t$ are the map frame 2D position, $\theta$ represents the odometry frame \chris{@Xiangyu, there is only odometry and please don't use IMU frame.}to the map frame rotation. $v$ is the scale ratio of the estimated velocity from IMU to the filter velocity. Whenever relative position updates $\mathbf{I_{t}} = {[I_{tx},I_{ty}]^T}$ in odometry frame is received, the states are propagated as followed:
\begin{equation}
\begin{aligned}
\mathbf{s}_{t}=\mathbf{s}_{t-1}+M(\mathbf{I_{t}} \times (1+\phi(n_p))) + [0,0,\phi(n_r),0]^T 
\end{aligned}
\end{equation}
\begin{equation}
\begin{aligned}
M = {\left[\begin{array}{l l l}{v_t\cos\theta_{t-1}}&{-v_t\sin\theta_{t-1}}\\ {v_t\sin\theta_{t-1}}&{v_t\cos\theta_{t-1}}\\ {0}&{0}\\ {0}&{0}\end{array}\right]}
\end{aligned}
\end{equation}


where $\phi(n)$ represents random Gaussian noise with zero mean and variance $n$.


\begin{equation}
\begin{aligned}
\check{\mathbf{P}}_{t}=\mathbf{F}_{t-1}\check{\mathbf{P}}_{t-1}\mathbf{F}_{t-1}^{T}+\mathbf{N}_{k-1}\mathbf{Q}_{k-1}\mathbf{N}_{k-1}^{T}
\end{aligned}
\end{equation}
where $\check{\mathbf{P}}_{t}$ is the state covariance at time $t$, $Q_{t-1}$ is the noise covariance at time $t$. ${F}_{t-1}$ is the motion model jacobian with respect to the last state while ${N}_{t-1}$ is the motion model jacobian with respect to the noise of the last state. 

\subsubsection{Radio map correction stage} The weight $w$ of every particle is updated using its current location and the
likelihood map during this stage. To obtain a likelihood map $\mathbf{f_{map}}$ in the preprocessing phase, the crowdsourced radio map is fed into the Gaussian kernel density estimation algorithm \cite{} and then normalized.  The importance weight of every particle is updated using its current location states.
Like a regular particle filter algorithm, the ratio of effective particles is monitored during the process, and if it is lower than a certain threshold, the standard systematic method \cite{} is used for the particle resampling.

\subsubsection{WiFi Kalman correction}

 In the proposed system, the measurement model relates the current state $\mathbf{\check{s}}_t$ to WiFi localization at time t $\mathbf{L_{t}}$$ = [x_t, y_t]^T$ in the map frame is defined as:
\begin{equation}
\begin{aligned}
 \mathbf{L_{t}}=\left[\begin{array}{c c c c}{{1}}&{{0}}&{{0}}&{{0}}\\ {{0}}&{{1}}&{{0}}&{{0}}\end{array}\right]\mathrm{s}_{t}+{\bf noise}_{k}\;
\end{aligned}
\end{equation}

Then incorporate covariances to compute the Kalman Gain:
\begin{equation}
\begin{aligned}
\mathbf{K}_t = \mathbf{\check{P}}_t \mathbf{H}_t^T \left(\mathbf{H}_t \mathbf{\check{P}}_t \mathbf{H}_t^T + \mathbf{M}_t \mathbf{R}_t \mathbf{M}_t^T \right)^{-1} 
\end{aligned}
\end{equation}
where $R_t$ is the measurement noise covariance where we used our estimated uncertainty from the WiFi model, $\mathbf{H}_t$ is the measurement model jacobian with respect to the last state, $\mathbf{R}_t$ is the measurement model jacobian with respect to the noise. The Kalman Gain $K_t$ is used to correct the predicted state:
\begin{equation}
\begin{aligned}
\mathbf{\hat{s}}_t = \mathbf{\check{s}}_t + \mathbf{K}_t \left(\mathbf{L}_t - [x_t,y_t]^T\right)
\end{aligned}
\end{equation}
Covariance correction is done through:
\begin{equation}
\begin{aligned}
\mathbf{\hat{P}}_t = \left(\mathbf{Identity} - \mathbf{K}_t \mathbf{H}_t \right)\mathbf{\check{P}}_t
\end{aligned}
\end{equation}

